% GENERATED FILE. Edit canonical lessons, then run npm run generate:books.
% canonical-source-hash: fnv1a64:ec5267a6bd200ce5
% canonical-lessons: BN-C269-uchit, BN-W269-letter-u, BN-C269-uchit-te-hobe, BN-C269-tai-sejonno, BN-R269-advice-and-result

\chapter{Ought To, Have To, and So}
\label{ch:bn-ought-to-have-to-and-so}
\hlchaptermodality{269}

\begin{chapteropening}
\textbf{By the end of this chapter:} \emph{I can say what someone ought to do with uchit and has to do with -te hôbe, draw a result with tāi and sejônyô, and read and write \bn{উ} in \bn{উচিত}.}

Complete BN-R269-advice-and-result: give advice, a duty and a result.
\end{chapteropening}

\section[\texorpdfstring{\bn{উচিত} (uchit)}{উচিত (uchit)}]{\bn{উচিত} (uchit) — ought to; proper}

\label{lesson:BN-C269-uchit}



\begin{quote}
Before the new one: say \textquotedblleft{}because I'm tired\textquotedblright{} and \textquotedblleft{}because of the storm\textquotedblright{}. (\emph{kārôn āmi klāntô}, \emph{jhôṛer jônyô}.)
\end{quote}

\subsection*{You'll want to know: \bn{উচিত}}
\textbf{\bn{উচিত}} — \emph{uchit} — \textquotedblleft{}ought to; proper\textquotedblright{}.

\begin{quote}
\textbf{\bn{তোমার} \bn{ডাক্তারের} \bn{কাছে} \bn{যাওয়া} \bn{উচিত।}}
\end{quote}

\begin{quote}
\emph{tomār ḍāktārer kāchhe jāoyā uchit.} — \textquotedblleft{}You ought to go to the doctor.\textquotedblright{}
\end{quote}

The one who ought to is in the \textbf{-\bn{র}} form, \textbf{\bn{তোমার}}, as in \textbf{\bn{তোমার} \bn{মনে} \bn{হয়}}; the action is the dictionary form, \textbf{\bn{যাওয়া}}; \textbf{\bn{উচিত}} comes last. Word for word: \textquotedblleft{}your going to the doctor is proper\textquotedblright{}.

\begin{cousinweb}
\textbf{\bn{উচিত}} (\emph{uchit}) is a tatsama: Sanskrit \emph{ucita}, which Monier-Williams gives as \textquotedblleft{}customary\textquotedblright{} and \textquotedblleft{}proper, suitable\textquotedblright{}. CDIAL has no entry for it.
\end{cousinweb}

\subsection*{Your turn}
Say these aloud:

\begin{itemize}
\raggedright
  \item \emph{uchit}
  \item \textquotedblleft{}You ought to go to the doctor\textquotedblright{} — \emph{tomār ḍāktārer kāchhe jāoyā uchit}
  \item \textquotedblleft{}You ought to rest\textquotedblright{} — \emph{tomār bishrām neoyā uchit}
\end{itemize}

\begin{tcolorbox}[breakable,colback=teal!4,colframe=teal!35!black,title={Before you move on}]
What does \textbf{\bn{উচিত}} mean? (\textquotedblleft{}Ought to; proper\textquotedblright{}.) Who goes in the \textbf{-\bn{র}} form? (The one who ought to: \textbf{\bn{তোমার}}.)
\end{tcolorbox}
\section[\texorpdfstring{\bn{উ} (u)}{উ (u)}]{\bn{উ} (u) — the letter at the start of \bn{উচিত}}

\label{lesson:BN-W269-letter-u}



\begin{quote}
Before the new one: say \textbf{\bn{উচিত}}. Then read \textbf{\bn{তু}} (\emph{tu}): the sign \textbf{\bn{ু}} under \textbf{\bn{ত}} is the vowel of this lesson.
\end{quote}

\subsection*{Script you'll notice: \bn{উ}}
\textbf{\bn{উ}} — \emph{u}. After a consonant you write it as the sign \textbf{\bn{ু}}: \textbf{\bn{তুমি}}. At the start of a word, with nothing to ride on, it is the letter \textbf{\bn{উ}}: \textbf{\bn{উচিত}}. \textbf{\bn{আ}} and \textbf{\bn{া}} are the same pair.

In the print of this book, \textbf{\bn{উ}} is \textbf{\bn{ড}} with the raised hook of \textbf{\bn{ট}} on top. Below the headline it has \textbf{\bn{ড}}'s body: a slanting left arm that is the left side of a round bowl, a short stem hanging from the headline in the middle, and a curve joining that stem to the bowl's right side, which ends in a point. Above the headline it has \textbf{\bn{ট}}'s hook: a curve standing on the headline, its free end a tip at the upper left.

\subsection*{Writing: \bn{উ} — copy the printed shape}
Put your pen beside \textbf{\bn{উ}} and copy the shape you can see, slowly and large. Check the copy against the print: the hook above the headline, the left arm and bowl, the short middle stem and its curve, the point on the right.

\begin{quote}
This book does not yet tell you where to start \bn{উ} or which way the pen goes. That is not written down here until it can be written down with a source. Copying the printed shape needs no such source, and it is how the shape becomes yours in the meantime.
\end{quote}

\subsection*{Your turn}
\begin{itemize}
\raggedright
  \item \emph{Find:} \bn{উ} at the front of \textbf{\bn{উচিত}}
  \item \emph{Write it:} \bn{ড}, then \bn{ট}, then \bn{উ}, from the print
  \item \emph{Write it:} \textbf{\bn{উচিত}}
  \item \emph{Write it:} \textbf{\bn{কারণ}}, from memory
  \item \emph{Write it:} \textbf{\bn{উঠল}} (\emph{uṭhlô}), the past of \textbf{\bn{ওঠা}}, now that you write \textbf{\bn{উ}}
\end{itemize}

\begin{tcolorbox}[breakable,colback=teal!4,colframe=teal!35!black,title={Before you move on}]
Which letter opens \textbf{\bn{উচিত}}? (\textbf{\bn{উ}}.) Which two letters does it put together? (\textbf{\bn{ড}} below the headline, \textbf{\bn{ট}}'s hook above.) Where does the sign \textbf{\bn{ু}} go? (Under its consonant.)
\end{tcolorbox}
\section[\texorpdfstring{… \bn{উচিত} / -\bn{তে} \bn{হবে} (… uchit / -te hôbe)}{… উচিত / -তে হবে (… uchit / -te hôbe)}]{… \bn{উচিত} / -\bn{তে} \bn{হবে} (uchit / -te hôbe) — ought to, have to}

\label{lesson:BN-C269-uchit-te-hobe}



\begin{quote}
Say \textquotedblleft{}You ought to go to the doctor\textquotedblright{}. (\emph{tomār ḍāktārer kāchhe jāoyā uchit}.)
\end{quote}

\subsection*{You'll want to know: ought to, ought not}
\textbf{\bn{উচিত}} gives advice. \textbf{\bn{না}} after it gives a warning:

\begin{quote}
\textbf{\bn{রাতে} \bn{তোমার} \bn{একা} \bn{যাওয়া} \bn{উচিত} \bn{না।}}
\end{quote}

\begin{quote}
\emph{rāte tomār ekā jāoyā uchit nā.} — \textquotedblleft{}You shouldn't go alone at night.\textquotedblright{}
\end{quote}

For yourself, \textbf{\bn{আমার}}: \textbf{\bn{আমার} \bn{বিশ্রাম} \bn{নেওয়া} \bn{উচিত}} (\emph{āmār bishrām neoyā uchit}), \textquotedblleft{}I ought to rest\textquotedblright{}.

\subsection*{You'll want to know: have to}
What you must do, not only ought to, is \textbf{-\bn{তে} \bn{হবে}}: the one who must in the \textbf{-\bn{কে}} form, the verb in its \textbf{-\bn{তে}} form, then \textbf{\bn{হবে}}.

\begin{quote}
\textbf{\bn{কাল} \bn{আমাকে} \bn{স্টেশনে} \bn{যেতে} \bn{হবে।}}
\end{quote}

\begin{quote}
\emph{kāl āmāke sṭeshône jete hôbe.} — \textquotedblleft{}Tomorrow I have to go to the station.\textquotedblright{}
\end{quote}

The \textbf{-\bn{তে}} form: \textbf{\bn{যাওয়া}} $\to$ \textbf{\bn{যেতে}} (\emph{jete}), \textbf{\bn{খাওয়া}} $\to$ \textbf{\bn{খেতে}} (\emph{khete}), \textbf{\bn{করা}} $\to$ \textbf{\bn{করতে}} (\emph{kôrte}), \textbf{\bn{আসা}} $\to$ \textbf{\bn{আসতে}} (\emph{āste}). In the past, \textbf{\bn{হবে}} becomes \textbf{\bn{হল}}: \textbf{\bn{আমাকে} \bn{হেঁটে} \bn{যেতে} \bn{হল}} (\emph{āmāke hẽṭe jete hôlô}), \textquotedblleft{}I had to walk\textquotedblright{}.

\subsection*{Your turn}
Say these aloud:

\begin{itemize}
\raggedright
  \item \textquotedblleft{}You shouldn't go alone at night\textquotedblright{} — \emph{rāte tomār ekā jāoyā uchit nā}
  \item \textquotedblleft{}Tomorrow I have to go to the station\textquotedblright{} — \emph{kāl āmāke sṭeshône jete hôbe}
  \item \textquotedblleft{}I had to walk\textquotedblright{} — \emph{āmāke hẽṭe jete hôlô}
\end{itemize}

\begin{tcolorbox}[breakable,colback=teal!4,colframe=teal!35!black,title={Before you move on}]
\textquotedblleft{}Ought not\textquotedblright{}? (\textbf{\bn{উচিত} \bn{না}}.) Which form of \textquotedblleft{}I\textquotedblright{} goes with \textbf{-\bn{তে} \bn{হবে}}? (\textbf{\bn{আমাকে}}.) \textbf{\bn{যাওয়া}} in the \textbf{-\bn{তে}} form? (\textbf{\bn{যেতে}}.)
\end{tcolorbox}
\section[\texorpdfstring{\bn{তাই} / \bn{সেজন্য} (tāi / sejônyô)}{তাই / সেজন্য (tāi / sejônyô)}]{\bn{তাই} / \bn{সেজন্য} (tāi / sejônyô) — so, that is why}

\label{lesson:BN-C269-tai-sejonno}



\begin{quote}
Say \textquotedblleft{}I had to walk\textquotedblright{}. (\emph{āmāke hẽṭe jete hôlô}.) Then \textbf{\bn{তাই}}, \textquotedblleft{}so\textquotedblright{}, from chapter one hundred and fifty-nine.
\end{quote}

\subsection*{You'll want to know: so}
\textbf{\bn{কারণ}} puts the result first and the reason after it. \textbf{\bn{তাই}} turns that round: the reason first, then \textquotedblleft{}so\textquotedblright{}, then the result.

\begin{quote}
\textbf{\bn{বাইরে} \bn{ঝড়} \bn{হচ্ছিল}, \bn{তাই} \bn{আমরা} \bn{বাড়িতে} \bn{থাকলাম।}}
\end{quote}

\begin{quote}
\emph{bāire jhôṛ hôchchhilô, tāi āmrā bāṛite thāklām.} — \textquotedblleft{}There was a storm outside, so we stayed at home.\textquotedblright{}
\end{quote}

\subsection*{You'll want to know: that is why}
\textbf{\bn{সেজন্য}} (\emph{sejônyô}), \textquotedblleft{}for that, that is why\textquotedblright{}, is \textbf{\bn{সে}}, \textquotedblleft{}that\textquotedblright{}, and the \textbf{\bn{জন্য}} of the last chapter. It usually opens a new sentence:

\begin{quote}
\textbf{\bn{আমার} \bn{জ্বর} \bn{ছিল।} \bn{সেজন্য} \bn{আমাকে} \bn{বাড়িতে} \bn{থাকতে} \bn{হল।}}
\end{quote}

\begin{quote}
\emph{āmār jbôr chhilô. sejônyô āmāke bāṛite thākte hôlô.} — \textquotedblleft{}I had a fever. That's why I had to stay at home.\textquotedblright{}
\end{quote}

\subsection*{Your turn}
Say these aloud:

\begin{itemize}
\raggedright
  \item \textquotedblleft{}There was a storm, so we stayed at home\textquotedblright{} — \emph{jhôṛ hôchchhilô, tāi āmrā bāṛite thāklām}
  \item \textquotedblleft{}I had a fever. That's why I had to stay at home.\textquotedblright{} — \emph{āmār jbôr chhilô. sejônyô āmāke bāṛite thākte hôlô.}
  \item \textquotedblleft{}The train was late, so I took a bus\textquotedblright{} — \emph{ṭren derite elô, tāi āmi bāse gelām}
\end{itemize}

\begin{tcolorbox}[breakable,colback=teal!4,colframe=teal!35!black,title={Before you move on}]
Which comes first with \textbf{\bn{তাই}}, the reason or the result? (The reason.) And with \textbf{\bn{কারণ}}? (The result.) What is \textbf{\bn{সেজন্য}} made of? (\textbf{\bn{সে}} and \textbf{\bn{জন্য}}: \textquotedblleft{}for that\textquotedblright{}.)
\end{tcolorbox}
\section[(review)]{Review — ought to, have to, and so}

\label{lesson:BN-R269-advice-and-result}



\begin{quote}
Say \textquotedblleft{}You ought to rest\textquotedblright{} and \textquotedblleft{}so we stayed at home\textquotedblright{}. (\emph{tomār bishrām neoyā uchit}; \emph{tāi āmrā bāṛite thāklām}.)
\end{quote}

\subsection*{Your turn}
\begin{itemize}
\raggedright
  \item \emph{Say it:} I think you ought to go to the doctor.
  \item \emph{Say it:} You shouldn't go alone at night.
  \item \emph{Say it:} Tomorrow I have to go to school early.
  \item \emph{Say it:} The bus didn't come, so we walked.
  \item \emph{Say it:} Mira was tired. That's why she stayed at home.
  \item \emph{Write it:} \textbf{\bn{উচিত}}, then \textbf{\bn{কারণ}}
\end{itemize}

\begin{tcolorbox}[breakable,colback=teal!4,colframe=teal!35!black,title={Before you move on}]
Advice or a must: \textbf{\bn{উচিত}} or \textbf{-\bn{তে} \bn{হবে}}? (Advice is \textbf{\bn{উচিত}}.) Which opens the result: \textbf{\bn{কারণ}} or \textbf{\bn{তাই}}? (\textbf{\bn{তাই}}.) Which letter opens \textbf{\bn{উচিত}}? (\textbf{\bn{উ}}.)
\end{tcolorbox}
